\clearpage
\fancyhead[L]{\small SONNENDURCHMESSER.}
\hypertarget{native-530-lower}{}\label{native-530-lower}
\begin{center}\small Göttingische gelehrte Anzeigen. Stück 95. S. 937..940. 16. Juni 1810.\\[3pt]\rule{.72\linewidth}{.3pt}\end{center}

Im 113. Stück dieser \emph{Göttingischen gelehrten Anzeigen} vom vorigen Jahr [S. 517 dieses Bandes] haben wir die bemerkenswerten Ergebnisse mitgeteilt, die Herr \textsc{von Lindenau} aus den Beobachtungen in Greenwich von 1750... 1786 über den Sonnendurchmesser gewonnen hatte. Seither hat der Astronom seine Untersuchungen auch auf die neueren Beobachtungen von 1787... 1798 ausgedehnt und die Ergebnisse unserem Herrn Professor \textsc{Gauss} in einem Schreiben vom 23. Mai mitgeteilt. Es ist interessant, dass sie sehr gut mit den früheren übereinstimmen. Deshalb verdienen sie umso mehr volle Aufmerksamkeit, zumindest in praktischer Hinsicht. Zunächst zeigt sich auch hier wieder eine Art periodischer Ungleichheit der auf die mittlere Entfernung reduzierten Halbmesser je nach Jahreszeit, wie die folgende Übersicht zeigt.

\begin{center} \begin{tabular}{l|c||l|c} & Sonnenhalbmesser & & Sonnenhalbmesser \\ \hline Januar & $959''70$ & Juli & $959''22$ \\ Februar & $959.99$ & August & $959.98$ \\ März & $960.41$ & September & $960.19$ \\ April & $959.80$ & Oktober & $960.10$ \\ Mai & $959.81$ & November & $960.07$ \\ Juni & $959.00$ & Dezember & $958.75$ \\ \end{tabular} \end{center}

Diese Beobachtungen lassen sich recht gut darstellen, wenn man in der Formel an der angeführten Stelle auf S. 1125 [S. 519 dieses Bandes] (wo übrigens durch einen Druckfehler 0.9631 statt 0.09631 steht) $\frac{bb}{aa} - 1 = \frac{1}{92}$ setzt. Der Wert ist also etwas kleiner als der dort bestimmte, aber ebenfalls positiv.

Die mittleren Sonnenhalbmesser aus allen einzelnen Beobachtungen der verschiedenen Jahre gibt Herr \textsc{von Lindenau} wie folgt an:

\begin{center} \begin{tabular}{l|c|c} & \begin{tabular}{c}Sonnenhalb-\\messer\end{tabular} & \begin{tabular}{c}Anzahl der\\Beobachtungen\end{tabular} \\ \hline 1787 & $959''26$ & 62 \\ 1788 & $959.10$ & 75 \\ 1789 & $959.68$ & 95 \\ 1790 & $959.33$ & 80 \\ 1791 & $959.51$ & 91 \\ 1792 & $960.24$ & 65 \\ 1793 & $959.59$ & 75 \\ 1794 & $959.96$ & 62 \\ 1795 & $960.05$ & 64 \\ 1796 & $960.02$ & 48 \\ 1797 & $959.99$ & 70 \\ 1798 & $960.03$ & 72 \\ \end{tabular} \end{center}

\SourceLeaf{531}
Zur scheinbaren allmählichen Abnahme des Sonnenhalbmessers liefert die Zusammenstellung der Beobachtungen \textsc{Maskelyne}s aus dreiunddreißig Jahren in drei Zeitabschnitten folgende Ergebnisse:

\begin{center} \begin{tabular}{c|c} & Sonnenhalbmesser \\ \hline 1765... 1775 & $961''66$ \\ 1776... 1786 & $960.22$ \\ 1787... 1798 & $959.77$ \\ \end{tabular} \end{center}

Aus 116 von \textsc{Piazzi} in den Jahren 1791... 1793 beobachteten horizontalen Sonnenhalbmessern findet Herr \textsc{von Lindenau} $961''21$, aus den Beobachtungen \textsc{Bradley}s von 1755... 1760 ebenso $961''86$. Die enge Übereinstimmung dieser Messungen aus so verschiedenen Zeiten bestätigt, dass die Größe des Sonnenkörpers unverändert bleibt. Den Grund für die allmähliche Abnahme während des langen Zeitraums der Beobachtungen \textsc{Maskelyne}s sieht Herr \textsc{von Lindenau} dagegen in einer Veränderung der Sehkraft des Beobachters, bei dem die Irradiation mit zunehmendem Alter geringer geworden zu sein scheint. Herr \textsc{von Lindenau} bezeichnet dies selbst nur als eine Vermutung, die einer weiteren Untersuchung nicht unwürdig ist.

Der beträchtliche Unterschied zwischen den horizontalen und vertikalen Halbmessern findet sich auch in den neueren Beobachtungen wieder. Die Berechnungen des Herrn \textsc{von Lindenau} ergeben hierzu Folgendes:

\begin{center}\small\setlength{\tabcolsep}{3pt} \begin{tabular}{l|c|c||l|c|c} & \begin{tabular}{c}Vertikalhalb-\\messer\end{tabular} & \begin{tabular}{c}Anzahl der\\Beobachtungen\end{tabular} & & \begin{tabular}{c}Vertikalhalb-\\messer\end{tabular} & \begin{tabular}{c}Anzahl der\\Beobachtungen\end{tabular} \\ \hline 1787 & $962''42$ & 72 & 1793 & $963''22$ & 49 \\ 1788 & $962.44$ & 102 & 1794 & $963.13$ & 72 \\ 1789 & $963.10$ & 91 & 1795 & $963.22$ & 60 \\ 1790 & $963.03$ & 66 & 1796 & $962.40$ & 61 \\ 1791 & $963.10$ & 57 & 1797 & $962.25$ & 50 \\ 1792 & $963.30$ & 48 & 1798 & $962.10$ & 53 \\ \end{tabular} \end{center}

Das arithmetische Mittel aus allen diesen Werten ist $962''82$ und damit um $2''96$ größer als das Mittel aus den horizontalen Halbmessern. Aus den Beobachtungen von 1765... 1786 ergab sich für diesen Unterschied $2''80$. Aus 71 Beobachtungen \textsc{Piazzi}s findet Herr \textsc{von Lindenau} für den mittleren vertikalen Halbmesser $963''01$ und für den horizontalen einen um $1''80$ kleineren Wert. Auch Beobachtungen von \textsc{Bouvard} ergaben für Herrn \textsc{von Lindenau} einen ähnlichen Unterschied.

Diese Übereinstimmung der Ergebnisse aus den Beobachtungen dreier verschiedener Beobachter ist in der Tat sehr auffällig. Wenn jedoch eine länglich-ellipsoidische Gestalt des Sonnenkörpers aus anderen, überwiegenden Gründen kaum als annehmbar gelten kann, dürfte man eher geneigt sein, den Grund dieses Unterschieds anderswo zu suchen, und vielleicht am ehesten in der Art der Beobachtung. Denn bei der Messung der Zenitdistanzen bringt man gewöhnlich den Sonnenrand \emph{von außen} mit dem Horizontalfaden in Berührung. Daher muss man für den oberen Sonnenrand eine um die halbe Fadendicke zu kleine Zenitdistanz und für den unteren eine um denselben Betrag zu große finden. Die auf diese Art bestimmten Sonnendurchmesser müssen folglich um die ganze Fadendicke zu groß ausfallen.
